% Standalone problem document, generated from the adjacent README.md.
% Compile with XeLaTeX. Mathematical content is maintained in Markdown.
\documentclass[11pt,a4paper]{article}
\usepackage[fontsize=10.5pt]{scrextend}
\usepackage[margin=22mm,headheight=15pt,headsep=9mm,footskip=11mm]{geometry}
\usepackage{fontspec}
\setmainfont{texgyrepagella}[Extension=.otf,UprightFont=*-regular,BoldFont=*-bold,ItalicFont=*-italic,BoldItalicFont=*-bolditalic]
\setsansfont{texgyreheros}[Extension=.otf,UprightFont=*-regular,BoldFont=*-bold,ItalicFont=*-italic,BoldItalicFont=*-bolditalic]
\setmonofont{texgyrecursor}[Extension=.otf,UprightFont=*-regular,BoldFont=*-bold,ItalicFont=*-italic,BoldItalicFont=*-bolditalic,Scale=MatchLowercase]
\usepackage{amsmath,amssymb}
\usepackage{unicode-math}
\setmathfont{texgyrepagella-math.otf}
\usepackage{xcolor,microtype,enumitem,needspace,fancyhdr}
\usepackage{bookmark}
\definecolor{ink}{HTML}{17344B}
\definecolor{accent}{HTML}{197D80}
\definecolor{muted}{HTML}{596773}
\hypersetup{colorlinks=true,linkcolor=accent,urlcolor=accent,pdftitle={SP-14 - Widom's
canonical distribution conjecture for Toeplitz
eigenvalues},pdfauthor={Open Problems in Numerical Linear Algebra}}
\urlstyle{same}
\setlength{\parindent}{0pt}
\setlength{\parskip}{5pt plus 1pt minus 1pt}
\linespread{1.04}
\setlength{\emergencystretch}{3em}
\widowpenalty=10000
\clubpenalty=10000
\displaywidowpenalty=10000
\allowdisplaybreaks[1]
\setlist{nosep,leftmargin=1.5em}
\providecommand{\tightlist}{\setlength{\itemsep}{0pt}\setlength{\parskip}{0pt}}
\providecommand{\pandocbounded}[1]{#1}
\pagestyle{fancy}
\fancyhf{}
\fancyhead[L]{\sffamily\footnotesize\color{muted}OPEN PROBLEMS IN NUMERICAL LINEAR ALGEBRA}
\fancyhead[R]{\sffamily\bfseries\footnotesize\color{accent}SP-14}
\fancyfoot[L]{\parbox[b]{0.9\textwidth}{\sffamily\fontsize{8}{10}\selectfont\color{muted}Editorial ratings; a bounded search cannot certify that a problem remains unsolved.\\Literature check: 2026-09-11}}
\fancyfoot[R]{\sffamily\footnotesize\color{muted}\thepage}
\renewcommand{\headrulewidth}{0.3pt}
\renewcommand{\headrule}{\hbox to\headwidth{\color{accent}\leaders\hrule height \headrulewidth\hfill}}
\makeatletter
\renewcommand{\section}{\Needspace{5\baselineskip}\@startsection{section}{1}{0pt}{12pt plus 2pt minus 2pt}{4pt}{\sffamily\bfseries\large\color{ink}}}
\renewcommand{\subsection}{\Needspace{4\baselineskip}\@startsection{subsection}{2}{0pt}{9pt plus 1pt minus 1pt}{3pt}{\sffamily\bfseries\normalsize\color{ink}}}
\makeatother
\setcounter{secnumdepth}{0}
\begin{document}
{\sffamily\small\color{muted}Eigenvalues and inverse problems\par}
\vspace{3pt}
{\raggedright\hyphenpenalty=10000\exhyphenpenalty=10000\sffamily\bfseries\fontsize{21}{24}\selectfont\color{ink}Widom's
canonical distribution conjecture for Toeplitz eigenvalues\par}
\vspace{7pt}
{\sffamily\small\textbf{Difficulty:} extreme\quad\textbf{Importance:} interesting
to the community\par}
{\sffamily\small\bfseries\color{accent}Status: Partially resolved\par}
\vspace{4pt}
{\color{accent}\hrule height 0.8pt}
\vspace{6pt}
\textbf{Topic:} Nonnormal Toeplitz eigenvalue asymptotics

\textbf{Rating rationale:} A general theorem must distinguish symbols
whose Toeplitz spectra follow their values from the markedly different
behavior of analytic symbols. This is a longstanding obstacle in the
spectral analysis of nonnormal structured matrices.

\subsection{Statement}\label{statement}

Let \(\mathbb T=\{z\in\mathbb C:|z|=1\}\) and
\(a\in C(\mathbb T;\mathbb C)\). Define

\[
a_k=\frac1{2\pi}\int_0^{2\pi}a(e^{it})e^{-ikt}\,dt,
\qquad T_n(a)=(a_{j-k})_{j,k=0}^{n-1}.
\]

Say that \(a\) extends analytically to an inner annulus if, for some
\(0<r<1\), there is a holomorphic function on \(r<|z|<1\) extending
continuously to \(|z|=1\) with boundary value \(a\). Define extension to
an outer annulus analogously using \(1<|z|<R\) for some \(R>1\).

\textbf{Conjecture (Widom).} If \(a\) has neither such inner nor such
outer extension, then for every continuous compactly supported
\(F:\mathbb C\to\mathbb C\),

\[
\lim_{n\to\infty}\frac1n\sum_{j=1}^n F(\lambda_j(T_n(a)))
=\frac1{2\pi}\int_0^{2\pi}F(a(e^{it}))\,dt.
\]

Eigenvalues are counted with algebraic multiplicity. This conclusion is
called \emph{canonical eigenvalue distribution}. The conjecture asserts
that failure requires a one-sided analytic extension; it does not assert
failure for every symbol that has one.

\subsection{Known cases and numerical
significance}\label{known-cases-and-numerical-significance}

Canonical distribution holds for real-valued symbols and for the Tilli
class, whose essential range has empty interior and connected
complement. Widom also proved cases with nonsmooth symbols tracing
Jordan curves. The general continuous-symbol assertion remains the
target.

Toeplitz singular-value distributions are much better understood than
their nonnormal eigenvalue counterparts. This question determines when
sampling the symbol predicts the bulk eigenvalues, a basic issue for
structured eigensolvers. It is distinct from
\href{https://github.com/ajt60gaibb/OpenProblemsInNLA/blob/main/eigenvalues-and-inverse-problems/SP-06/README.md}{SP-06},
which asks for real spectra in every finite order, and from
individual-eigenvalue expansion problems.

\newpage
\subsection{References and status
check}\label{references-and-status-check}

\begin{itemize}
\tightlist
\item
  H. Widom, \emph{Eigenvalue distribution of nonselfadjoint Toeplitz
  matrices and the asymptotics of Toeplitz determinants in the case of
  nonvanishing index}, Operator Theory: Advances and Applications 48
  (1990), 387--421. Historical originating source, as identified by the
  explicit restatements below.
\item
  J. M. Bogoya, A. Böttcher and S. M. Grudsky, \emph{Asymptotics of
  individual eigenvalues of a class of large Hessenberg Toeplitz
  matrices}, OTAA 220 (2012), 77--95.
  \href{https://www.math.cinvestav.mx/~grudsky/Papers/116.pdf}{Author
  manuscript}, §1, p.2 after (1.1), states the annulus conjecture
  explicitly.
\item
  M. Bogoya, S.-E. Ekström, S. Serra-Capizzano and P. Vassalos,
  \emph{Matrix-less methods for the spectral approximation of large
  non-Hermitian Toeplitz matrices: A concise theoretical analysis and a
  numerical study}, Numerical Linear Algebra with Applications 31
  (2024), e2545. \href{https://doi.org/10.1002/nla.2545}{DOI};
  \href{https://uu.diva-portal.org/smash/get/diva2\%3A1824576/FULLTEXT01.pdf}{institutional
  PDF}. The introduction before Theorem 1 reaffirms the conjecture and
  states the Tilli-class result.
\end{itemize}

On 2026-09-11, checked the explicit source statements and searched
Widom, canonical distribution, analytic annuli, and nonselfadjoint
Toeplitz proof/counterexample combinations, including 2025--2026
results. No full resolution was found. Other conjectures bearing Widom's
name, about determinant or trace asymptotics, do not settle this
eigenvalue-distribution statement. This is a bounded status check.
\end{document}
